\section{Detailed mathematical derivations}
\label{app:derivations}

This appendix contains the calculations that are used in the article but whose
length would interrupt it: the element matrices of the Legendre basis, five
worked matrix elements, the construction of the blocks \eqref{eq:block} from the
weak equations, the edge-matching system, the normalization of the modal power,
and one end-to-end numerical evaluation of the reported share.

\subsection{Element matrices of the Legendre basis}

Three coordinates appear in the formulation and they must be separated before
any integral is written, because the orthonormality of the basis and the size of
the element matrices both depend on which one is meant.  The physical coordinate
is $z\in[-d/2,d/2]$.  The non-dimensional coordinate of the field equations is
$\zeta=2z/d\in[-1,1]$, with $\dd z=(d/2)\dd\zeta$.  The basis is generated on a
third, the half-interval variable
\begin{equation}
\label{eq:app_measures}
s=\frac{z}{d}=\frac{\zeta}{2}\in\Bigl[-\tfrac12,\tfrac12\Bigr],
\qquad
\dd s=\frac{\dd z}{d}=\frac{\dd\zeta}{2},
\end{equation}
and each element matrix in this appendix is an integral over $s$ with the
measure $\dd s$.  The basis is
\begin{equation}
\label{eq:app_basis}
b_j(s)=\sqrt{2j+1}\,P_j(2s),\qquad j=0,\dots,N-1,
\end{equation}
with $N$ the number of basis functions per field.  Since
$\int_{-1}^{1}P_iP_j\dd\zeta=2\delta_{ij}/(2j+1)$, the Gram matrix in the $s$
measure is
\begin{equation}
\label{eq:app_mass}
\int_{-1/2}^{1/2}b_ib_j\,\dd s
=\frac{2j+1}{2}\int_{-1}^{1}P_iP_j\,\dd\zeta=\delta_{ij},
\end{equation}
the first factor being the Jacobian of $s=\zeta/2$ and the second the Legendre
normalization.  The same integral in the $\zeta$ measure is
$\int_{-1}^{1}b_ib_j\dd\zeta=2\delta_{ij}$: the two differ by just the Jacobian, and the mass matrix that multiplies $\dot{\vect q}$ in the weak form is
the identity only in the $s$ measure.  Hence all statements below are made with $\dd s$, and the derivative in $S_{ij}$ is $\dd_s=\dd/\dd s$; in the $\zeta$
coordinate $\dd_s=2\dd/\dd\zeta$, which is the factor of two that appears in all strain-gradient entries.

The stiffness matrix is $S_{ij}=\int b_i'b_j'\dd s$.  For $i=j$ the chain rule
gives $b_j'(s)=2\sqrt{2j+1}\,P_j'(2s)$, so
\begin{equation}
\label{eq:app_S}
S_{jj}=4(2j+1)\int_{-1/2}^{1/2}\bigl(P_j'(2s)\bigr)^{2}\dd s
=2(2j+1)\,j(j+1),
\end{equation}
where $\int_{-1}^{1}(P_j')^{2}\dd\zeta=j(j+1)$ is the classical companion of the
Legendre equation and the factor $2$ collects $4\times\tfrac12$.  The weighted
form of that identity is the one that follows directly from the equation, and it is derived here because this step is usually skipped.  Multiplying
$\bigl[(1-\zeta^{2})P_j'\bigr]'=-j(j+1)P_j$ by $P_k$ and integrating by parts, the
boundary term carrying the factor $(1-\zeta^{2})$ vanishes at both ends and
\begin{equation}
\label{eq:app_S_weighted}
\int_{-1}^{1}(1-\zeta^{2})P_j'P_k'\,\dd\zeta
=\frac{2j(j+1)}{2j+1}\,\delta_{jk},
\end{equation}
which is the orthogonality of the derivatives under the weight
$(1-\zeta^{2})$; \eqref{eq:app_S} is its unweighted companion and is used here
only through the closed form it implies.  The discrete matrix checks both, and so
does an independent quadrature rule of a different order
(\cref{app:repro}): the assembled $S_{jj}$ equals $2j(j+1)(2j+1)$ for every
$j\le N-1$ at all degrees used in this article, to $2\times10^{-11}$ or better
in absolute value.

For $i\neq j$ the parity of $P_i'P_j'$ leaves $S_{ij}\neq0$ only when $i$ and $j$
have the same parity, and
\begin{equation}
\label{eq:app_S_off}
S_{ij}=2\sqrt{(2i+1)(2j+1)}\int_{-1}^{1}P_i'P_j'\,\dd\zeta ,
\end{equation}
which explains the structure of \eqref{eq:app_S_matrix}: the matrix is
real and symmetric and it vanishes between basis functions of opposite parity,
but the same-parity entries are not confined to a band.  At $N=6$ the entry
$S_{15}=22.978$ couples indices four apart, and at any $N$ the row $i$ has
non-zero entries in each column $j\le i$ of the same parity, so the matrix is
not tridiagonal and has no fixed bandwidth; what bounds it is the parity rule,
not a band.  At $N=6$ the matrix is
\begin{equation}
\label{eq:app_S_matrix}
\mathsf S=
\begin{pmatrix}
0&0&0&0&0&0\\
0&12&0&18.330&0&22.978\\
0&0&60&0&80.498&0\\
0&18.330&0&168&0&210.60\\
0&0&80.498&0&360&0\\
0&22.978&0&210.60&0&660
\end{pmatrix},
\end{equation}
whose diagonal is \eqref{eq:app_S} entry by entry.  Equation
\eqref{eq:app_S_matrix} is the entire gradient-energy content of the
discretization for a single field: it is real, symmetric and independent of the
material, and each strain-gradient block of \eqref{eq:block} is this matrix
multiplied by an elastic modulus and by $\lambda^{2}$.

Two further element matrices are needed.  The potential space of the grounded
problem is spanned by
\begin{equation}
\label{eq:app_pot}
\psi_j(s)=b_j(s)-r_j\,b_{j+2}(s),\qquad
r_j=\sqrt{\frac{2j+1}{2j+5}},\qquad j=0,\dots,N_\varphi-1 ,
\end{equation}
which vanishes at $s=\pm\tfrac12$ for every $j$ because $P_j(\pm1)=(\pm1)^j$ and
$r_j$ is exactly the ratio that makes the two terms cancel; $N_\varphi=N-2$,
since $\psi_j$ needs $b_{j+2}$ up to $j=N_\varphi-1$, i.e.\ $b_{N-1}$.  The two
matrices that couple the charge are then distinguished by which index is the test
index.  Testing Gauss's law with $\psi_i$ against the carrier basis $b_j$ gives
\begin{equation}
\label{eq:app_charge}
\bigl(\mathsf C\bigr)_{ij}\equiv\int_{-1/2}^{1/2}\psi_i\,b_j\,\dd s
=\delta_{ij}-r_i\,\delta_{j,i+2},
\qquad i=0,\dots,N_\varphi-1,\quad j=0,\dots,N-1 ,
\end{equation}
that is a unit matrix minus one band above the diagonal, with the coefficient
$r_i$ carrying the \emph{row} index.  The reciprocal matrix that appears in the
carrier equation is its transpose,
\begin{equation}
\label{eq:app_charge_T}
\bigl(\mathsf C^{\mathsf T}\bigr)_{ij}\equiv\int_{-1/2}^{1/2}b_i\,\psi_j\,\dd s
=\delta_{ij}-r_j\,\delta_{i,j+2},
\end{equation}
which has the band below the diagonal and the coefficient $r_j$ on the column
index.  Both appear in \eqref{eq:block}: $-\mathsf C$ in the potential row and
$-\mathsf C^{\mathsf T}$ in the carrier row, entering through the single term
$-\int W_i^{*}W_j$ whose test and trial functions share the potentials.  Writing
$r_j$ where $r_i$ belongs, or transposing the two, gives a matrix that differs
from the assembled one by $0.2981$ in the largest entry at degree four, so the orientation is stated here and not left to the reader.  Finally
the dielectric matrix is
$\kappa_{33}\int\psi_i'\psi_j'\dd s+\kappa_{11}k_y^{2}\int\psi_i\psi_j\dd s$,
with
\begin{equation}
\label{eq:app_diel}
\int_{-1/2}^{1/2}\psi_i'\psi_j'\,\dd s
=S_{ij}-r_jS_{i,j+2}-r_iS_{i+2,j}+r_ir_jS_{i+2,j+2},
\end{equation}
in which the second and third terms are equal because $\mathsf S$ is symmetric,
and the first term is not $S_{ij}$ alone: the grounded basis removes a band of
the stiffness exactly as it removes a band of the mass.

\subsection{Five worked matrix elements}

The entries below are evaluated at the operating point $k_xd=0.7362$,
$k_yd=0.1298$, $\omega=7.6334\,c_s/d$ and at polynomial degree $p=4$, i.e.\
$N=p+1=5$ basis functions per field, $n_{\mathrm{u}}=3N=15$, $n_\varphi=N-2=3$,
$n_\mu=N=5$ and $n_{\mathrm{w}}=5p+3=23$ unknowns.  That is the smallest degree at
which the parity structure and the grounded bands are both visible; the reported
results use $p=22$, with $p=18$ to $34$ examined, where the same formulas hold with larger matrices.
All entries are in the scaled units of the solver: stresses divided by
$c_{44}$, electric displacement by $\sqrt{c_{44}\kappa_{33}}$ and electric field by $\sqrt{c_{44}/\kappa_{33}}$, permittivities by
$\kappa_{33}$, so the numbers are dimensionless.

\begin{enumerate}
\item \emph{Compressibility (mass) entry}, carrier row and column:
\begin{equation}
\label{eq:app_w1}
-\int_{-1/2}^{1/2}b_ib_j\,\dd s=-\delta_{ij}
\qquad\text{so that}\qquad
\bigl(\mathsf A_{00}\bigr)_{\mu_0\mu_0}=-1 ,
\end{equation}
reproduced as $-1$ to $7\times10^{-15}$.
\item \emph{Transport entry.}  The carrier gradient enters the weak form without a
power of $k_x$ through its transverse part and with one power through its axial
part, so the $k_x^{0}$ block holds the transverse wavenumber and the $k_x^{2}$
block the axial one:
\begin{equation}
\label{eq:app_w2}
\bigl(\mathsf A_{00}\bigr)_{\mu_i\mu_j}\Big|_{\text{transport}}
=-\frac{\ii\alpha_{\mathrm d}}{\omega}\Bigl[k_y^{2}\,\delta_{ij}+S_{ij}\Bigr],
\qquad
\bigl(\mathsf A_{11}\bigr)_{\mu_i\mu_j}\Big|_{\text{transport}}
=-\frac{\ii\alpha_{\mathrm d}}{\omega}\,\delta_{ij} ,
\end{equation}
so that the polynomial operator reproduces
$-\ii\alpha_{\mathrm d}\bigl(k_x^{2}+k_y^{2}+S_{ij}\bigr)/\omega$.  With
$\alpha_{\mathrm d}/\omega=0.108509$ and $k_y^{2}=0.016848$ the
diagonal of \eqref{eq:app_w2} is
\begin{equation}
\label{eq:app_w2_values}
-\ii\,1.828\times10^{-3},\quad -\ii\,1.3039,\quad -\ii\,6.5124,\quad
-\ii\,18.2314 \qquad (j=0,\dots,3),
\end{equation}
for $S_{jj}=0,12,60,168$, and the $k_x^{2}$ block is
$-\ii\,0.108509$.  The two pieces together reproduce the assembled coefficient of
the polynomial to $3.6\times10^{-15}$.  The growth with $j$ is the discrete
statement that a higher Legendre mode stores more gradient energy, and hence that diffusion penalizes thickness structure.
\item \emph{Carrier--potential entry.}  With the orientations fixed by
\eqref{eq:app_charge}--\eqref{eq:app_charge_T}, the potential row of
$\mathsf A_{00}$ carries $-\mathsf C$ and the carrier row
$-\mathsf C^{\mathsf T}$:
\begin{equation}
\label{eq:app_w3}
\bigl(\mathsf A_{00}\bigr)_{\varphi_i\mu_j}
=-Q\bigl(\delta_{ij}-r_i\,\delta_{j,i+2}\bigr),
\qquad
\bigl(\mathsf A_{00}\bigr)_{\mu_i\varphi_j}
=\bigl(\mathsf A_{00}\bigr)_{\varphi_j\mu_i},
\end{equation}
so the diagonal is $-Q=-0.074635$ and the band is $+Qr_i$, that is $+0.033378$,
$+0.048860$, $+0.055630$ and $+0.059538$ for $i=0,\dots,3$; the first of these is
$Q/\sqrt{5}$ because $r_0=1/\sqrt{5}$, and the sign is positive because the
band of \eqref{eq:app_charge} is negative and the block has the overall minus
of $-\int W_i^{*}W_j$.  In the carrier row the same numbers appear at
$(i,i-2)$, below the diagonal, with the coefficient $r_{i-2}$ carrying the
\emph{trial} (potential) index; the two matrices are transposes of one another, and
the assembled blocks reproduce \eqref{eq:app_w3} to $9\times10^{-16}$.
\item \emph{Gradient-stiffness entry}, out-of-plane row and column:
\begin{equation}
\label{eq:app_w4}
\lambda^{2}\frac{c_{33}}{c_{44}}\,S_{ij}
=0.0625\times5.1395\times S_{ij}
\;=\;3.85,\ 19.27,\ 53.97,\ 115.64 \quad (j=1,2,3,4),
\end{equation}
growing as $j^{3}$ because $S_{jj}\propto j^{3}$: this is the growth that makes
the discrete operator badly scaled before row equilibration, and $115.64$ at $j=4$ is the largest of these entries at
this degree.
\item \emph{Inertia entry}:
\begin{equation}
\label{eq:app_w5}
-\omega^{2}\delta_{ij}=-58.27\,\delta_{ij}\qquad(\omega=7.6334).
\end{equation}
The ratio of the inertia entry to the $j=1$ gradient entry is
$58.27/3.85=15.1$.  It compares two matrix entries at one fixed frequency and
nothing more: the frequency of a thickness mode is an eigenvalue of the coupled
problem \eqref{eq:quartic}, not the ratio of two entries, and no mode frequency is
inferred from it here.
\end{enumerate}

\subsection{Construction of the blocks from the weak equations}

Each of the nine blocks $\mathsf A_{rs}$ is the sum of the eight bilinear forms
listed in \eqref{eq:block}.  Written out with the operators of
\eqref{eq:sector_blocks} and the integrals of the first subsection of this appendix,
the block is
\begin{equation}
\label{eq:app_block}
\begin{aligned}
\mathsf A_{rs}={}&\underbrace{\bigl\langle\mathsf B_r,\mathsf C\mathsf B_s
\bigr\rangle}_{\text{elastic}}
+\underbrace{\lambda^{2}\bigl\langle\partial_\zeta\mathsf B_r,
\mathsf C\,\partial_\zeta\mathsf B_s\bigr\rangle}_{\text{gradient}}
+\underbrace{\bigl\langle\mathsf G_r,\mathsf e\mathsf B_s\bigr\rangle
+\bigl\langle\mathsf B_r,\mathsf e^{\mathsf T}\mathsf G_s\bigr\rangle}_{\text{piezoelectric}}\\
&-\underbrace{\bigl\langle\mathsf G_r,\boldsymbol\kappa\mathsf G_s\bigr\rangle
}_{\text{dielectric}}
+\underbrace{f\bigl[\bigl\langle\mathsf G_r^{(2)},\partial_\zeta\mathsf B_s^{(2)}
\bigr\rangle+\bigl\langle\partial_\zeta\mathsf B_r^{(2)},\mathsf G_s^{(2)}
\bigr\rangle\bigr]}_{\text{flexoelectric}}
-\underbrace{\frac{\ii\alpha_{\mathrm d}}{\omega}
\bigl\langle\mathsf M\mathsf G_r,\mathsf M\mathsf G_s\bigr\rangle}_{\text{transport}}\\
&-\underbrace{\bigl\langle\mathsf W_r,\mathsf W_s\bigr\rangle}_{\text{carrier}}
-\underbrace{\omega^{2}\bigl\langle\mathsf U_r,\mathsf U\bigr\rangle\delta_{r0}
\delta_{s0}}_{\text{inertia}},
\end{aligned}
\end{equation}
with $\langle\vect a,\vect b\rangle=\int_{-1/2}^{1/2}\vect a^{\mathsf H}\vect b\,
\dd s$ in the measure of \eqref{eq:app_measures}.

The powers of $k_x$ are counted by the block indices $(r,s)$ and by nothing else,
and it is written out here because double counting is the usual error.
The strain of a single mode is split once and for all as
\begin{equation}
\label{eq:app_strain_split}
\varepsilon_{ij}=\varepsilon^{(0)}_{ij}+\ii k_x\,\varepsilon^{(1)}_{ij},
\end{equation}
with both coefficient vectors independent of $k_x$:
\begin{center}
\renewcommand{\arraystretch}{1.1}
\begin{tabular}{@{}lcc@{}}
\toprule
component & $\varepsilon^{(0)}_{ij}$ & $\varepsilon^{(1)}_{ij}$ \\
\midrule
$\varepsilon_{xx}$ & $0$ & $b_j$ in the $\widehat u_1$ slot \\
$\varepsilon_{yy}$ & $\ii k_yb_j$ in the $\widehat u_2$ slot & $0$ \\
$\varepsilon_{zz}$ & $b_j'$ in the $\widehat u_3$ slot & $0$ \\
$\varepsilon_{yz}$ & $b_j'$ in $\widehat u_2$, $\ii k_yb_j$ in $\widehat u_3$ & $0$ \\
$\varepsilon_{xz}$ & $b_j'$ in the $\widehat u_1$ slot & $b_j$ in the $\widehat u_3$ slot \\
$\varepsilon_{xy}$ & $\ii k_yb_j$ in the $\widehat u_1$ slot & $b_j$ in the $\widehat u_2$ slot \\
\bottomrule
\end{tabular}
\end{center}
no component contains more than one power of $\ii k_x$, so the split terminates at
first order.  The strain gradient is split the same way, with $\partial_x$ acting
on the test factor and on the trial factor separately, and the block is defined as
the coefficient of $(-\ii k_x)^{r}(\ii k_x)^{s}$ in \eqref{eq:app_block}:
\begin{equation}
\label{eq:app_ar_s}
\mathsf A_{rs}=\bigl\langle\varepsilon^{(r)},\mathsf C\varepsilon^{(s)}\bigr\rangle
+\lambda^{2}\bigl\langle\partial_z\varepsilon^{(r)},\mathsf C\partial_z\varepsilon^{(s)}\bigr\rangle
+\cdots ,
\end{equation}
the remaining terms of \eqref{eq:app_block} being split in the same way.  Two consequences follow in this language.  First, the leading coefficient
is the gradient energy of the axial part of the strain alone,
\begin{equation}
\label{eq:app_A22}
\mathsf A_{22}=\lambda^{2}\bigl\langle\varepsilon^{(1)},\mathsf C\varepsilon^{(1)}\bigr\rangle
\neq0 ,
\end{equation}
which is where the quartic term and the four wave families come from; its rank is $n_{\mathrm u}$ and not $n_{\mathrm w}$, because $\varepsilon^{(1)}$ is
supported on the displacement slots only.  Secondly, no strain component contains
two powers of $\ii k_x$, which is why $\mathsf A_{02}=\mathsf A_{20}=0$ and the polynomial is a quartic in the $(r,s)$ bookkeeping of \eqref{eq:polynomial}, not a sextic: the
$k_x^{4}$ term needs one factor from each side of the gradient energy and cannot
be produced by a single factor carrying two.

The same bookkeeping fixes where the flexoelectric constant can appear.  The only
non-zero component of $f_{ijkl}$ used in this article is $f_{3333}$, which pairs
the field $E_z$ with the gradient component $\eta_{zzz}=\partial_z\varepsilon_{zz}$;
neither contains a power of $k_x$, since $E_z=-\partial_z\varphi$ and
$\varepsilon_{zz}$ is axial, so $f$ enters $\mathsf A_{00}$ and no other block.
That is a statement about the operator, and it explains why a non-zero $f$ can
change the elastic-like term and the double traction without touching the quartic
coefficient $\mathsf A_{22}$: the two constants sit in different blocks, and
neither of them is a $k_x$ power.

\subsection{The edge-matching system}

At the edge $x=0$ the conditions \eqref{eq:bc_edge} are imposed weakly, one row
block per condition and each on the test functions of the field it multiplies.
The normal there is the $x$ direction, not the $z$ direction: the generalized
traction $\bar\sigma_{1j}$ on the $n_{\mathrm u}=3N$ displacement test functions,
the double traction $\tau_{1j1}\equiv\tau_{jxx}$ on the same $n_{\mathrm u}$ test
functions, the normal electric displacement $D_x$ on the $n_\varphi$ potential
test functions and the normal carrier flux $J_x$ on the $n_\mu$ carrier test
functions.  The face quantities $\bar\sigma_{iz}$, $\tau_{izz}$ and $D_z$ of
\eqref{eq:bc_faces_mech}--\eqref{eq:bc_grounded} belong to $z=\pm d/2$; they are
the natural conditions of the expansion in $z$ and are not the edge conditions.
The number of rows is thus
\begin{equation}
\label{eq:app_edge_rows}
m=2n_{\mathrm u}+n_\varphi+n_\mu ,
\end{equation}
with the double-traction block present whenever the hyperstress is non-zero.  At the edge only the axial strain gradient remains:
\begin{equation}
\label{eq:app_edge_tau}
\tau_{ijx}=\ell^{2}c_{ijkl}\eta_{klx} .
\end{equation}
The two gradient couplings of \eqref{eq:tauZnO} act on the normal-derivative
component: $f$ pairs $E_z$ with $\eta_{zzz}$ and $H$ pairs $\Delta n$ with the
same component, so neither appears in \eqref{eq:app_edge_tau}.  The block exists if and only if $\ell\neq0$, independently of $f$ and of $H$: it
is present in cases A, B and D below, and it is empty in the classical limit
$\ell=f=H=0$, where $\tau_{ijx}\equiv0$ and the double-traction statement
degenerates to $0=0$ and does not give a boundary condition.  The case $\ell=0$ with
$f\neq0$ retains a hyperstress but has no axial trace to pair it with, which is
why it is not a case of this model.  Between those cases the edge problem changes
order, and with it the state:

\begin{center}
\renewcommand{\arraystretch}{1.15}
\begin{tabular}{@{}lllccl@{}}
\toprule
case & parameters & $\tau_{ijx}$ & $N_s$ & $m$ & double traction \\
\midrule
A & $\ell\neq0$, $f\neq0$, $H\neq0$ & $\ell^{2}c_{ijkl}\eta_{klx}$ & $2n_{\mathrm w}+2n_{\mathrm u}$ & $2n_{\mathrm u}+n_\varphi+n_\mu$ & required \\
B & $f=0$, $\ell\neq0$ & $\ell^{2}c_{ijkl}\eta_{klx}$ & $2n_{\mathrm w}+2n_{\mathrm u}$ & $2n_{\mathrm u}+n_\varphi+n_\mu$ & required \\
C & $\ell=0$, $f\neq0$ & $0$ (no axial gradient) & $2n_{\mathrm w}$ & $n_{\mathrm u}+n_\varphi+n_\mu$ & \emph{not available} \\
D & $H=0$, $\ell\neq0$ & $\ell^{2}c_{ijkl}\eta_{klx}$ & $2n_{\mathrm w}+2n_{\mathrm u}$ & $2n_{\mathrm u}+n_\varphi+n_\mu$ & required \\
E & $\ell=f=H=0$ & $0$ & $2n_{\mathrm w}$ & $n_{\mathrm u}+n_\varphi+n_\mu$ & absent (identically zero) \\
\bottomrule
\end{tabular}
\end{center}

\noindent Case C is not a case of the present model: the hyperstress would be non-zero while the
second-order state of case E, which has no $\partial_x u$ trace to pair the
double traction with, is the only one available at $\ell=0$.  Thus the formulation admits $\ell=0$ only in the classical limit $f=H=0$, and states that restriction explicitly; it does not assume that $\ell=0$ implies a second-order model.  No
number in this article depends on case C: every reported run has $\ell\neq0$.

With the edge ansatz \eqref{eq:edge_ansatz} and the modal state
$\vect y^{(j)}\in\mathbb{C}^{N_s}$ (the $n_{\mathrm w}$ displacement and
potential coefficients, the $n_{\mathrm u}$ axial traces, and the $n_{\mathrm w}$
and $n_{\mathrm u}$ conjugate rows), the system is
\begin{equation}
\label{eq:app_edge}
\vect y(0)=\vect y_{\text{inc}}+\sum_{j=1}^{M}a_j\,\vect y^{(j)},
\qquad
\mathsf B\vect a=-\vect b_{\text{inc}},
\end{equation}
where the boundary operator acts on a \emph{state} and the boundary matrix on the
modal index,
\begin{equation}
\label{eq:app_edge_R}
\mathsf R=\begin{pmatrix}
\mathsf R_{\Sigma}\\
\mathsf R_{\tau}\\
\mathsf R_{D}\\
\mathsf R_{J}
\end{pmatrix}\in\mathbb{C}^{m\times N_s},
\qquad
\mathsf B=\bigl[\mathsf R\vect y^{(1)}\ \cdots\ \mathsf R\vect y^{(M)}\bigr]
\in\mathbb{C}^{m\times M},
\qquad
\vect b_{\text{inc}}=\mathsf R\,\vect y_{\text{inc}}\in\mathbb{C}^{m},
\end{equation}
so the three objects have three different second dimensions and only the last two
agree.  They do agree, identically, not by accident: the number of
decaying states is half the state dimension, $M=N_s/2=n_{\mathrm w}+n_{\mathrm u}$
in cases A, B and D, and $M=n_{\mathrm w}$ in case E where $N_s=2n_{\mathrm w}$;
with $n_{\mathrm w}=n_{\mathrm u}+n_\varphi+n_\mu$ from \eqref{eq:nw},
\begin{equation}
\label{eq:app_edge_square}
M=n_{\mathrm w}+n_{\mathrm u}=2n_{\mathrm u}+n_\varphi+n_\mu=m
\quad(\ell\neq0),
\qquad
M=n_{\mathrm w}=n_{\mathrm u}+n_\varphi+n_\mu=m
\quad(\ell=0),
\end{equation}
so the edge system is square in all cases and never over- or under-determined.
At degree four ($n_{\mathrm u}=15$, $n_\varphi=3$, $n_\mu=5$, $n_{\mathrm w}=23$)
this is $N_s=76$, $m=M=38$; at degree eighteen ($57$, $17$, $19$, $93$) it is
$N_s=300$, $m=M=150$; at the reported degree twenty-two ($69$, $21$, $23$, $113$)
it is $N_s=364$, $m=M=182$.  The singular-value decomposition is still used, but as a square solve instead of a least-squares fit: the requirement is that the
rank equal the column count at a cut-off of $10^{-12}$ of the largest singular
value, and a deficient system is reported as such and is not regularized.  Row and
column scaling are applied as before, and the conditioning at the operating point
is $1.4\times10^{2}$ after scaling.
\subsection{Normalization of the modal power}

The power \eqref{eq:Px_int} is bilinear in a mode and its conjugate, so the power
carried by the $j$-th outgoing mode scales with the square of its excitation
amplitude and with the normalization of its eigenvector.  The eigenvector is
normalized by $\vect y^{(j)}\to\vect y^{(j)}/\sqrt{\mathsf y^{\mathsf H}\mathsf
M\mathsf y}$, which fixes its mass norm but leaves the \emph{sign} of the power free:
the physical power is
\begin{equation}
\label{eq:app_power}
P_j=|a_j|^{2}\,\mathcal{P}_x\bigl(\vect y^{(j)}\bigr),
\end{equation}
and the degree-invariance of $P_j$ is the test that the normalization has been
applied correctly.  At the operating point, for the backward mode,
\begin{equation}
\label{eq:app_power_numbers}
\begin{aligned}
n=18:&\quad \mathcal{P}_x=-0.0178923,\quad |a|=1.03848,\quad P_B=-0.0192958,\\
n=22:&\quad \mathcal{P}_x=-0.0154248,\quad |a|=1.11846,\quad P_B=-0.0192957,\\
n=26:&\quad \mathcal{P}_x=-0.0132455,\quad |a|=1.20697,\quad P_B=-0.0192957,\\
n=30:&\quad \mathcal{P}_x=-0.0113777,\quad |a|=1.30227,\quad P_B=-0.0192956,
\end{aligned}
\end{equation}
in which the raw power falls by \SI{36}{\percent} and the squared amplitude rises by
\SI{57}{\percent} across the four degrees while the physical power is constant to
the six digits shown.  That constancy is the
numerical statement of \eqref{eq:app_power} and it is why the reported
share is degree-independent at the $10^{-6}$ level of
\cref{tab:refinement}.

\subsection{One end-to-end evaluation of the reported share}

The reported number at the operating point is assembled from five quantities,
each of which is stored with each row of the data: the raw backward power
$\mathcal{P}_x$, the scattering amplitude $|a_B|$, the incident power at the
reference plane, the two attenuation exponents and the reference distance
$L=10\,d$.  The five are quoted below from \emph{one} row, the reported degree
$n=22$, and not mixed with the neighboring degrees of
\eqref{eq:app_power_numbers}, because the raw power and the amplitude both change
with degree even though their product does not:
\begin{equation}
\label{eq:app_RB}
R_B=\frac{|a_B|^{2}\bigl|\mathcal{P}_x\bigr|\,e^{2\Imag k_{Bx}L}}
{P_{\text{inc}}\,e^{2\Imag k_{ix}L}}
=\frac{0.0192957\times0.9998954}{1.0001081}
=\frac{0.0192937}{1.0001081}=0.0192916 ,
\end{equation}
with, at $n=22$,
\begin{equation}
\label{eq:app_RB_inputs}
|a_B|=1.1184598,\quad
|a_B|^{2}=1.250952,\quad
\mathcal{P}_x=-0.015424813,\quad
|a_B|^{2}\bigl|\mathcal{P}_x\bigr|=0.0192957 ,
\end{equation}
the last product being the power the backward port actually carries at the edge.
The factor $e^{2\Imag k_{Bx}L}=0.9998954$ brings it to $x=-L$ with
$\Imag k_{Bx}=-5.2306\times10^{-6}$ per unit $d$, and the incident mode carries
$e^{2\Imag k_{ix}L}=1.0001081$ with $\Imag k_{ix}=+5.4027\times10^{-6}$; both
exponents are degree-independent to the seven digits quoted, so the degree of the
row enters \eqref{eq:app_RB} only through $|a_B|^{2}|\mathcal{P}_x|$.  What the
number measures is the ratio of two physical powers,
$-P_B/P_{\text{inc}}=0.0192916$, and the two exponential factors state at which
plane each power was taken.

The value and its meaning are two different statements and are kept apart.
Equation \eqref{eq:app_RB} evaluates the share at the \emph{operating point},
$1.93$ percent of the incident power at $x=-10\,d$.  The local maximum of the
same quantity over the frequency window, $R_B=0.9428718$, is a property of the
peak row, not of this one, and \eqref{eq:app_RB} neither computes nor implies it.  Reading the operating-point row along the plate, the backward mode carries
$1.93$ percent at $x=-10\,d$ and $1.93\,\%\times e^{-2.13\times10^{-4}}$ at
$x=-20\,d$: the decay of the share is governed by the \emph{difference} of the two
attenuation constants, not by either one of them.
